\scanpage{199}\mspage{148}
En effet, un morphisme dans $\mathrm{Cat}$ qui est $W$-lisse et à fibres $W$-asphériques est dans $W$ (et $=$ \inserted{c'est une $W$-équivalence} universelle \deleted{\unsure{ment}}\ldots)

\uline{Corollaire 2} \inserted{On suppose que $W$ satisfait $W(7\text{bis})$ (axiome du carré cocartésien).}\margin{\unsure{inutile}} \deleted{Considérons un}\inserted{Supposons le} carré $Q$ cocartésien dans $\mathrm{Cat}$${}^{*}$, et que $\alpha_1$ soit une immersion ouverte (resp. fermée) et $f_1$ lisse (resp. propre). Alors $Q$ est \inserted{$W$-}homotopiquement cocartésien.

% slanted
\margin{NB On n'utilise \unsure{pas|\deleted{pas}} vraiment que $Q$ soit cocartésien, plutôt que $Q$ soit \uline{cartésien}, et $f_2$ (donc aussi $\alpha_1$) une imm. ouverte (resp. fermée), et que pour $y\in Y\setminus X_2$, $X_{1y}$ soit $W$-asphérique. (Dans le cas où $Q$ est cocartésien $X_{1y}$ est $\simeq e$.)}

\margin{On a besoin ici de $W(7\text{bis})$ \uline{forme directe}, i.e. de $\mathrm{Loc}(6'\mathrm a)$ (notation de \uline{XVI})}

Il suffit d'établir le résultat non respé, l'autre en résulte par dualité (puisque $Q$ est $W$-\deleted{homotopiquement}\inserted{cocartésien} sss $Q^\circ$ l'est).

Par hypothèse $f_1$ est \inserted{$W$-}lisse, d'autre part $\alpha_1$ et $f_2$ sont des imm. ouvertes donc $W$-lisses, donc $f_1\alpha_1$ est \inserted{$W$-}lisse aussi. (D'ailleurs $\alpha_2$ aussi, car le carré est aussi cartésien donc $\alpha_2$ déduit de $f_1$ \inserted{(qui est $W$-lisse par hypothèse)} par chgt. de base $X_2\to Y$.) Donc $\unsure{\varphi|f}_Q$ est lisse, et il suffit de voir que ses fibres sont $W$-asphériques. Le calcul \unsure{les} donne, pour $y$ dans $Y\setminus X_2$

\diagram{\(\unsure{\mathscr X|\mathcal X}_y\simeq\int\deleted{\unsure{\varnothing}} \begin{tikzcd}[ampersand replacement=\&,row sep=0pt,column sep=1.2em,cramped] \& e \\ \varnothing\arrow[ur]\arrow[dr] \& \\ \& \varnothing \end{tikzcd}\simeq e,\quad\text{OK.}\)}

Et pour $y$ dans $X_2$

\diagram{\(\mathscr X_{\unsure{y|\smash{\overset{\circ}{y}}}}\simeq\int \begin{tikzcd}[ampersand replacement=\&,row sep=0pt,column sep=1.2em,cramped] \& X_{1y} \\ X_{0y}\arrow[ur,"\sim"]\arrow[dr] \& \\ \& e \end{tikzcd},\)}
% archive sheet number: 199
